% ---------------------------------------------------------------
% FIG 1 — system pipeline (native TikZ, full width)
% ---------------------------------------------------------------
\begin{figure*}[t]
  \centering
  \begin{tikzpicture}[
    font=\small,
    >={Stealth[length=2.2mm,width=1.8mm]},
    blk/.style={rectangle, rounded corners=2.5pt, draw, line width=0.7pt,
      minimum height=12mm, text width=19mm, align=center, inner sep=3pt},
    flow/.style={-{Stealth[length=2.2mm,width=1.8mm]}, line width=1pt, draw=cGray},
    lbl/.style={font=\scriptsize\itshape, text=cGray, inner sep=1.5pt},
    node distance=8mm and 8.5mm]

    \node[blk, draw=cBlue,  fill=cBlueL]  (fw)   {DS-TWR firmware\\ \footnotesize DW1000\,+\,ESP32};
    \node[blk, draw=cAmber, fill=cAmberL, right=of fw]   (cal)  {Range calib.\\ \footnotesize delay\,+\,power};
    \node[blk, draw=cAmber, fill=cAmberL, right=of cal]  (nlos) {NLOS weight\\ \footnotesize first-path pwr};
    \node[blk, draw=cRed,   fill=cRedL,   right=of nlos] (fix)  {Robust\\ weighted LS};
    \node[blk, draw=cGreen, fill=cGreenL, right=of fix]  (est)  {MHE\,+\,rigid\\ \footnotesize body ($+$IMU)};
    \node[blk, draw=cGreen!60!black, fill=cGreen!22, right=of est, text width=22mm]
      (pose) {Pose\\ \footnotesize $x,y,z,\phi,\theta,\psi$};

    \node[blk, draw=cGray, fill=cGrayL, below=9mm of nlos, text width=24mm]
      (surv) {Anchor self-survey \\ \footnotesize inter-anchor ranging};

    \draw[flow] (fw)   -- node[lbl,above]{ranges,} node[lbl,below]{FP/RX pwr}  (cal);
    \draw[flow] (cal)  -- (nlos);
    \draw[flow] (nlos) -- (fix);
    \draw[flow] (fix)  -- (est);
    \draw[flow] (est)  -- (pose);
    \draw[flow] (surv.east) -| (fix.south);

    \begin{scope}[on background layer]
      \node[draw=cBlue!50, dashed, rounded corners, inner sep=5pt,
        fit=(fw), label={[cBlue,font=\scriptsize\bfseries]above:Radio (ESP32)}] {};
      \node[draw=cGreen!45, dashed, rounded corners, inner sep=5pt,
        fit=(cal)(pose)(surv),
        label={[cGreen!60!black,font=\scriptsize\bfseries]above:MATLAB host}] {};
    \end{scope}
  \end{tikzpicture}
  \caption{The DUNE pipeline. Five DW1000 anchors and one or two ESP32
  tags perform double-sided two-way ranging; the host calibrates each
  range for antenna delay and received-power bias, weights it by a
  first-path NLOS score, solves a robust least-squares fix, and fuses
  the stream in a moving-horizon estimator with a two-tag rigid-body
  constraint. Anchor coordinates are supplied by an inter-anchor
  self-survey, and roll, pitch, and height are read from the on-board
  IMU to complete the six-degree-of-freedom pose.}
  \label{fig:pipeline}
\end{figure*}

\section{Introduction}
\label{sec:intro}

\IEEEPARstart{A}{ccurate} indoor localization is a prerequisite for
autonomous mobile robots operating where satellite navigation is
unavailable. Among the candidate technologies, ultra-wideband (UWB)
ranging offers a favourable combination of centimetre-level time-of-flight
resolution, robustness to multipath relative to narrowband radio, and
low unit cost, and it has become a common choice for real-time
localization systems~\cite{alarifi2016uwb}. A single DW1000 transceiver
can time a radio round trip to a few centimetres, and the underlying
two-way ranging and multilateration are well established~\cite{dstwr_neirynck}.

In practice, however, a substantial gap separates this nominal ranging
resolution from the pose accuracy a robot experiences during operation.
Raw transceiver ranges carry hardware- and signal-dependent biases of
several decimetres; obstructed, non-line-of-sight (NLOS) paths inflate
the measured distance without any external indication; and the anchor
coordinates that anchor the whole solution must themselves be obtained,
conventionally by manual survey or an external camera on which the
deployment then depends. Closing this gap is largely a matter of
engineering that is seldom reported in reproducible detail, so each new
deployment tends to repeat it.

This paper documents a complete, low-cost UWB positioning stack, DUNE,
that addresses these issues end to end and is released in full. The
system localizes a small Ackermann rover from five fixed DW1000 anchors
and one or two ESP32-hosted tags, with all estimation performed on a
host computer (Fig.~\ref{fig:pipeline}). The design targets
reproducibility on commodity hardware rather than a single algorithmic
novelty, and its accuracy derives from the disciplined combination of
calibration, NLOS-aware estimation, and automatic infrastructure
localization described in the following sections.

The main contributions are:
\begin{enumerate}
  \item A complete and openly released pipeline from raw double-sided
  two-way ranging to a filtered vehicle pose, including a two-step radio
  calibration (antenna delay and received-power range bias). Unlike
  common open DW1000 libraries, which target a fixed small anchor set,
  the pipeline supports an arbitrary number of anchors ($N\!\ge\!4$) and
  degrades gracefully when only a subset responds.
  \item A continuous first-path-power NLOS weighting integrated with a
  robust weighted least-squares solver and a moving-horizon estimator
  that enforces a two-tag rigid-body constraint to recover heading.
  \item An automatic anchor self-localization procedure that recovers
  the anchor geometry from inter-anchor ranging alone, eliminating the
  need for an external survey of the infrastructure.
\end{enumerate}
